\documentclass{article}

% Language setting
% Replace `english' with e.g. `spanish' to change the document language
\usepackage[english]{babel}

% Set page size and margins
% Replace `letterpaper' with`a4paper' for UK/EU standard size
\usepackage[letterpaper,top=2cm,bottom=2cm,left=3cm,right=3cm,marginparwidth=1.75cm]{geometry}

% Useful packages
\usepackage{amsmath}
\usepackage{graphicx}
\usepackage[colorlinks=true, allcolors=blue]{hyperref}
\usepackage{indentfirst}
\usepackage{enumitem}
\usepackage{physics}
\usepackage{gensymb}
\usepackage{accsupp}

\title{Exercício 5  -  Cálculo Numérico para Geociência}


\begin{document}
\maketitle




Nesta atividade buscaremos retornar aos conceitos apresentados durante o curso, para que possamos revisar os conceitos vistos de zero de função.

Para isso, traremos o mesmo problema, de perfil térmico da crosta terrestre, que pode ser dado pela seguinte expressão:
\begin{align}
    \label{eq: EDO_crosta}
    \frac{d^2T(z)}{d^2z} = -\frac{H(z)}{k}.
\end{align}
Sendo $T$ a temperatura medida em uma profundidade $z$, $H$ é a geração de calor radiogênico e $k$ é a condutividade da rocha. Para esse problema, assumiremos que $k$ é uma constante e que o $H$ tem uma dependência da profundidade em que estamos, podendo ser escrito da seguinte forma: $H(z) = H_0e^{-z/h_r}$, sendo $h_r$ um parâmetro de escala e $H_0$ a produção de calor radiogênico.

Podemos resolver de forma que obtemos a seguinte expressão para a temperatura em função da profundidade
\begin{align}
    T(z) = T_0 + \qty(\frac{q_0}{k})z + \frac{H_0h_r^2}{k}(1-e^{-z/h_r}),
\end{align}
sendo $q_0$ o fluxo de calor com a superfície e $T_0$ a temperatura da superfície. Substituindo, com valores mais próximos dos reais obtidos na literatura, temos a seguinte representação.
\begin{align}
    T(z) = 10 + 16,8z + 79,10(1-e^{-0.1z}),
\end{align}


Com esse perfil de temperatura, temos agora como descobrir a partir de qual profundidade a alta temperatura começa a alterar de forma significativa a dureza das rochas, tornando-as mais dúcteis, interferindo diretamente nos nossos modelos sísmicos. Uma temperatura em que começa a deformação estrutural das rochas é de $500 \degree C$. Definido nosso problema, execute as seguintes tarefas:
\vspace{0.2 cm}
\hline 
\vspace{0.2 cm}
1) Implemente o método da bisseção (ou método da dicotomia) até a quarta iteração. Tome o intervalo inicial entre $[20,30]$. Escreva os resultados das iterações com $4$ algarismos significativos.

2) Dado o intervalo do item anterior, quantas iterações seriam necessárias para alcançarmos uma precisão menor que $0.1 \degree C$?

3) Implemente o método de Newton até a quarta iteração. Tome o intervalo inicial $[20,30]$. Escreva os resultados das iterações com $4$ algarismos significativos.

\vspace{2cm}
\noindent\centering\large{ESCREVA TODOS OS PASSOS COMO APRESENTADO DURANTE O CURSO}


\end{document}